\documentclass{article}

\usepackage{ctex}
\usepackage{amsmath}
\usepackage{amsfonts}
\usepackage{graphicx}
\usepackage{setspace}
\usepackage{amssymb}
\usepackage{latexsym}
\usepackage{amsthm, euscript, makeidx, color, mathrsfs}
\usepackage[numbers,sort&compress]{natbib}
\usepackage{hyperref}
\hypersetup{
    colorlinks=true,
    linkcolor=blue,
    urlcolor=cyan,
    citecolor=cyan
}

\usepackage{geometry}
\geometry{
	a4paper,
	left=2.5cm,
	right=2.5cm,
	top = 3cm,
	bottom = 2cm
}

\newtheorem {theorem}{Theorem}[section]
\newtheorem {lemma}[theorem]{{\bf Lemma}}
\newtheorem {corollary}[theorem]{{\bf Corollary}}
\newtheorem {proposition}[theorem]{{\bf Proposition}}
\theoremstyle{remark}
\newtheorem {remark}{{\bf Remark}}[section]
\newtheorem {conclusion}[remark]{Conclusion}
\newtheorem {condition}{{\bf Condition}}[section]
\theoremstyle{definition}
\newtheorem {definition}{{\bf Definition}}[section]
\newtheorem {example}{{\bf Example}}[section]
\theoremstyle{plain}
\numberwithin {equation}{section}
\newtheorem{assumption}{Assumption}
\numberwithin{assumption}{section}

\allowdisplaybreaks

\begin{document}

\title{机器学习概述}
\author{数据科学与大数据技术专业}
\date{讲义1}
\maketitle


\section{机器学习核心定义与分类}

\begin{enumerate}
  \item 为什么要用机器学习
    \begin{itemize}
      \item 很难手工编码出一个解决方案（例如视觉、语音）
      \item 系统需要适应不断变化的环境（例如垃圾邮件检测）
      \item 希望系统的表现优于人类程序员
      \item 隐私/公平（例如对搜索结果进行排序）
    \end{itemize}
  \item 机器学习与统计学的差别
    \begin{itemize}
      \item 它和统计学很相似…… 两个领域都试图从数据中发现模式。
      \item 两个领域都大量依赖微积分、概率论和线性代数，并共享许多核心算法。
      \item 但它不是统计学！
      \item 统计学更关注帮助人得出可靠结论；机器学习更关注构建自主智能体。
      \item 统计学更强调可解释性和数学严谨性；机器学习更强调预测性能、可扩展性和自主性。
    \end{itemize}
  \item 机器学习与AI的关系
    \begin{itemize}
      \item 如今，机器学习常常与人工智能一起被提及。
      \item AI 并不总是意味着一个基于学习的系统。
      \item 符号推理、基于规则的系统、树搜索等等。
    \end{itemize}
  \item 机器学习与人类学习的关系
    \begin{itemize}
      \item 人类学习是：非常高效地利用数据；一个完整的多任务系统（视觉、语言、运动控制等）；至少需要几年时间。
      \item 为了实现特定目的，机器学习最终并不一定要看起来像人类学习。
      \item 它可能会借鉴生物系统的想法，例如神经网络。
      \item 它的表现可能比人类更好，也可能更差。
    \end{itemize}
  \item 机器学习的历史
  \item 核心定义
    \begin{itemize}
      \item 机器学习是通过从数据中自动学习规律（如函数关系或概率分布），以对未知数据进行预测或决策的过程。
    \end{itemize}
  \item 主要任务类型
    \begin{itemize}
      \item 监督学习：利用带标签数据训练模型，预测或分类新样本。
        \begin{itemize}
          \item 回归任务：预测连续值（如房价、股价等）。
          \item 分类任务：判别离散标签（如手写数字识别、垃圾邮件分类）。
        \end{itemize}
      \item 无监督学习：利用无标签数据发现数据结构或特征。
        \begin{itemize}
          \item 聚类：将数据分组（如客户细分、图像分割）。
          \item 降维：将高维数据映射到低维空间，便于可视化和后续分析（如主成分分析PCA）。
        \end{itemize}
      \item 强化学习：通过与环境交互，学习最优策略以最大化长期收益（如自动驾驶、围棋AI）。
    \end{itemize}
\end{enumerate}

\textbf{问题思考}：

\begin{itemize}
  \item 监督学习和无监督学习的主要区别是什么？
  \item 请举例说明回归与分类任务的实际应用场景。
  \item 强化学习与前两者有何本质不同？
\end{itemize}


\section{机器学习四要素}

\begin{enumerate}
  \item 数据
    \begin{itemize}
      \item 训练集：用于模型参数的学习。
      \item 验证集：用于模型选择和调参，评估泛化能力。
      \item 测试集：最终评估模型性能（实际应用中常单独划分）。
    \end{itemize}
  \item 模型
    \begin{itemize}
      \item 线性模型：假设输入与输出之间存在线性关系（如线性回归、逻辑回归）。
      \item 非线性模型：通过基函数、神经网络等方式增强表达能力，适应复杂关系。
    \end{itemize}
  \item 学习准则
    \begin{itemize}
      \item 损失函数：
        衡量模型预测 $f(x)$ 与真实值 $y$ 之间的误差，记为 $\mathcal{L}(y, f(x))$（如均方误差、交叉熵）。
      \item 期望风险（Expected Risk）：
        模型在真实数据分布上的平均损失 $R(f) = \mathbb{E}[\mathcal{L}(y, f(x))]$。
        由于真实数据分布通常是未知的，因此期望风险很难直接计算。
      \item 经验风险（Empirical Risk）：
        模型在训练数据集上的平均损失 $\hat{R}_D(f) = \frac{1}{N} \sum_{i=1}^N \mathcal{L}(y_i, f(x_i))$。
      \item 经验风险最小化（ERM）：
        用经验风险来近似期望风险，并通过最小化经验风险来寻找最优模型。
        \begin{itemize}
            \item 问题：ERM 可能会导致过拟合，即模型在训练数据上表现很好，但在未知数据上表现很差。
        \end{itemize}
    \end{itemize}
  \item 优化算法
    \begin{itemize}
      \item 梯度下降法：通过计算损失函数对参数的梯度，迭代更新参数。
      \item 随机梯度下降（SGD）及其变体：每次用部分样本更新参数，提高效率，适合大规模数据。
      \item 学习率：控制每次参数更新的步长，影响收敛速度与稳定性。
    \end{itemize}
\end{enumerate}

\textbf{问题思考}：

\begin{itemize}
  \item 为什么需要划分训练集、验证集和测试集？
  \item 损失函数的选择对模型训练有何影响？
  \item 梯度下降法与随机梯度下降有何优缺点？
\end{itemize}


\section{泛化与正则化}

\begin{enumerate}
  \item 过拟合问题
    \begin{itemize}
      \item 过拟合：模型在训练集上表现良好，但在新数据上效果较差，常因模型过于复杂或数据不足。
    \end{itemize}
  \item 正则化方法
    \begin{itemize}
      \item 模型复杂度控制：
        \begin{itemize}
          \item L1正则化：鼓励参数稀疏，部分参数变为零。
          \item L2正则化：抑制参数过大，提升模型泛化能力。
          \item 数据增强：通过扩充训练样本提升模型鲁棒性。
        \end{itemize}
      \item 优化过程干预：
        \begin{itemize}
          \item 提前停止：基于验证集性能决定训练终止，防止过拟合。
          \item 权重衰减：在优化过程中对参数施加惩罚。
          \item Dropout：训练时随机失活部分神经元，减少模型对特定路径的依赖。
        \end{itemize}
    \end{itemize}
  \item 验证集作用
    \begin{itemize}
      \item 验证集用于评估模型在未见数据上的泛化性能，指导模型选择和训练停止时机。
    \end{itemize}
\end{enumerate}

\textbf{问题思考}：

\begin{itemize}
  \item 过拟合产生的原因有哪些？
  \item L1与L2正则化的区别和适用场景？
  \item 如何利用验证集提升模型泛化能力？
\end{itemize}





（注：详细定义和更多的内容请参考课程教材第一、二章）

\end{document}